\documentclass[10pt,a4paper]{amsart}
\usepackage[margin=19mm]{geometry}
\usepackage{amsmath,amssymb,mathtools}
\usepackage[scr=rsfs,cal=euler]{mathalfa}
\usepackage{xfrac}
\usepackage[unicode,hidelinks,bookmarks=false]{hyperref}
\newcommand{\ie}{i.e.\ }
\newcommand{\cf}{cf.\ }
\newcommand{\resp}{respectively\ }
\newcommand{\Subsection}{\subsection}
\providecommand{\nobreakdash}{\nobreak\hbox{-}}
\setlength{\emergencystretch}{2em}
\multlinegap0pt
\usepackage{amssymb}
\usepackage{mathtools}
\usepackage{tikz}
\usepackage{booktabs}
\usepackage[shortlabels]{enumitem}
\usepackage{caption}
\newtheorem{lemma}{Lemma}
\newtheorem{theorem}{Theorem}
\newtheorem{remark}{Remark}
\newcommand{\bg}{\bar{g}}
\newcommand{\momdens}{\boldsymbol{j}}
\newcommand{\tM}{\tilde{M}}
\newcommand{\tg}{\tilde{g}}
\newcommand{\tr}{\mathrm{tr}}
\title{The Angular Momentum Penrose Inequality}
\author{Da Xu}
\date{Source: arXiv:2512.06918v4 (CC BY 4.0)}
\begin{document}
\maketitle

\noindent\emph{Source and licence.} The Angular Momentum Penrose Inequality. Version arXiv:2512.06918v4 (CC BY 4.0), distributed by arXiv under the Creative Commons Attribution 4.0 International licence, http://creativecommons.org/licenses/by/4.0/, as recorded in the arXiv licence metadata for this exact version and shown on the abstract page https://arxiv.org/abs/2512.06918v4. Full licence notice: \texttt{licenses/arxiv-2512.06918-CC-BY-4.0.txt}.\par\smallskip

\noindent\textit{Editorial scope.} Counted unit: the article's theorem thm:J-conserve (source0.tex lines 2045--2058). The article's complete original proof of it is delivered with it (source0.tex lines 2068--2284, one \texttt{proof} environment of 217 lines, which the article places away from the statement). No mathematics is written, repaired or re-derived: the statement and the proof are the article's own lines, in the article's own notation, and no retained line is substituted or re-flowed.  Retained uncounted as the source-local context the proof needs: lines 1850--1856.  Nothing was omitted from any retained span, so the package carries no omission record.  No retained line contains a citation, so the wrapper carries no bibliography and no bibliography entry.  No retained line contains a figure environment, an image-inclusion command or drawing code, so no image is delivered and the package declares no asset dependency.  Editorial formatting only, in addition: an AMS \texttt{amsart} preamble that restates the article's own declarations and macros used by the retained lines, namely the environments \texttt{lemma}, \texttt{theorem}, \texttt{remark} on separate counters, together with the standard packages those lines need.  The retained environments print in this document as Lemma 1, Theorem 1, Remark 1 in order of appearance, whereas the article prints the same items with the numbering of its own full document.  The complete native source of the article as distributed in the arXiv e-print for this exact version is bundled unchanged as \texttt{sources/190530/source0.tex}.
\begin{lemma}[Level Set Homology Preservation]\label{lem:homology}
Let $u: \tM \to [0, 1]$ be the $p$-harmonic potential with $u|_\Sigma = 0$ and $u \to 1$ at infinity. For regular values $t_1, t_2 \in (0, 1)$, the level sets $\Sigma_{t_1}$ and $\Sigma_{t_2}$ are homologous in $M$:
\[
[\Sigma_{t_1}] = [\Sigma_{t_2}] \in H_2(M; \mathbb{Z}).
\]
In particular, all level sets are homologous to the outermost MOTS $\Sigma$.
\end{lemma}

\begin{theorem}[Angular Momentum Conservation]\label{thm:J-conserve}
Let $(M, g, K)$ be axisymmetric initial data with Killing field $\eta = \partial_\phi$, satisfying the \textbf{vacuum} constraint equations ($\mu = |\momdens| = 0$) in the exterior region $M_{\mathrm{ext}} := M \setminus \overline{\mathrm{Int}(\Sigma)}$. Let $u: \tM \to [0,1]$ be the axisymmetric $p$-harmonic potential with level sets $\Sigma_t = \{u = t\}$. Define the Komar angular momentum:
\[
J(t) := \frac{1}{8\pi}\int_{\Sigma_t} K(\eta, \nu_t) \, dA_t = \int_{\Sigma_t} \star_g \alpha_J,
\]
where $\alpha_J := \frac{1}{8\pi}K(\eta, \cdot)^\flat_g$ is the Komar 1-form and $\star_g$ is the Hodge star with respect to the physical metric $g$. Then:
\[
J(t) = J(0) = J \quad \text{for all } t \in [0, 1].
\]

\textbf{Key innovation:} The Komar angular momentum $J$ is a \textbf{topological invariant} under the vacuum hypothesis. By showing the Komar 1-form is co-closed ($d^\dagger\alpha_J = 0$) in vacuum, the flux integral becomes independent of the integration surface via Stokes' theorem. This cleverly circumvents the dynamical instability of angular momentum in general flows.

\textbf{Mechanism:} This conservation follows from de Rham cohomology, not dynamics. The vacuum momentum constraint implies the Komar 1-form is \textbf{co-closed}: $d^\dagger_g \alpha_J = 0$, equivalently $d(\star_g \alpha_J) = 0$. Since all level sets $\Sigma_t$ are homologous (Lemma~\ref{lem:homology}), Stokes' theorem implies the flux integral is independent of $t$.
\end{theorem}

\begin{proof}[Proof of Theorem~\ref{thm:J-conserve}]
The proof has three main components: (A) establishing that the Komar integral is metric-independent, (B) proving co-closedness $d^\dagger\alpha_J = 0$ for vacuum axisymmetric data, and (C) applying Stokes' theorem.

\textbf{Key Identity.} The central result is that for vacuum axisymmetric data ($\momdens_i = 0$ and $\mathcal{L}_\eta K = 0$), the Komar 1-form $\alpha_J = \frac{1}{8\pi}K(\eta, \cdot)^\flat$ satisfies:
\begin{equation}\label{eq:key-coclosed}
d^\dagger \alpha_J = -\star d\star \alpha_J = 0,
\end{equation}
which is equivalent to $d(\star_g \alpha_J) = 0$. This follows from the momentum constraint $\nabla^j K_{ij} = \nabla_i(\tr K) + 8\pi \momdens_i$ with $\momdens_i = 0$ (vacuum), combined with the Killing equation for $\eta$ (axisymmetry). Once \eqref{eq:key-coclosed} is established, Stokes' theorem immediately gives $J(\Sigma_{t_1}) = J(\Sigma_{t_2})$ for homologous surfaces.

\textbf{Part A: Metric-Independence of the Komar Integral.}
The Komar angular momentum is defined using the \textbf{physical} extrinsic curvature $K$ on $(M, g)$, while the AMO flow operates on $(\tM, \tg = \phi^4 \bg)$. We must show the conservation law transfers correctly, and that the Komar integral is independent of the choice of metric used to define the normal vector and area element.

\textbf{Definition of the Komar integral (metric-explicit).} The Komar 1-form is defined using the \textbf{physical} metric $g$:
\[
\alpha_J := \frac{1}{8\pi} K(\eta, \cdot)^\flat_g = \frac{1}{8\pi} K_{ij} \eta^i g^{jk} dx_k.
\]
This is a well-defined 1-form on the smooth manifold $M$, independent of any choice of metric for the integration surface.

For a 2-surface $\Sigma \subset M$, the Komar angular momentum is computed as follows. Let $\star \alpha_J$ denote the Hodge dual of $\alpha_J$ (a 2-form). Then:
\[
J(\Sigma) = \int_\Sigma \star \alpha_J.
\]
Alternatively, if we choose \textbf{any} Riemannian metric $\gamma$ on $M$ and let $\nu_\gamma$ be the $\gamma$-unit normal and $d\sigma_\gamma$ the $\gamma$-area element:
\[
J(\Sigma) = \int_\Sigma \alpha_J(\nu_\gamma) \, d\sigma_\gamma = \int_\Sigma K(\eta, \nu_\gamma) \cdot \frac{d\sigma_\gamma}{8\pi}.
\]

\textbf{Key claim: The integral is metric-independent.} Suppose $\gamma_1$ and $\gamma_2$ are two Riemannian metrics on $M$. We claim:
\[
\int_\Sigma \alpha_J(\nu_{\gamma_1}) \, d\sigma_{\gamma_1} = \int_\Sigma \alpha_J(\nu_{\gamma_2}) \, d\sigma_{\gamma_2}.
\]

\begin{proof}[Proof of metric-independence]
We prove this by showing both expressions equal the integral of a metric-independent 2-form.

\textit{Step (i): Construction of the flux 2-form.} Given the 1-form $\alpha_J$ on a 3-manifold $M$ and a 2-surface $\Sigma \subset M$, we construct the associated flux. Let $\iota: \Sigma \hookrightarrow M$ be the inclusion. Choose \textbf{any} smooth extension of the normal field: for any metric $\gamma$, extend $\nu_\gamma$ to a neighborhood $U \supset \Sigma$ as a vector field (still denoted $\nu_\gamma$).

Define the 2-form on $\Sigma$:
\[
\omega_\Sigma := \iota^*(\iota_{\nu_\gamma} \text{vol}_\gamma) \cdot \alpha_J(\nu_\gamma),
\]
where $\text{vol}_\gamma$ is the volume form of $\gamma$. We claim this is independent of $\gamma$.

\textit{Step (ii): Coordinate calculation.} Let $(y^1, y^2)$ be local coordinates on $\Sigma$ and extend to coordinates $(y^1, y^2, n)$ on $U$ where $n$ is a coordinate transverse to $\Sigma$ with $\Sigma = \{n = 0\}$. In these coordinates, the $\gamma$-unit normal is $\nu_\gamma = \frac{1}{|\partial_n|_\gamma}\partial_n + (\text{tangential corrections})$, the area element is $d\sigma_\gamma = |\partial_n|_\gamma \sqrt{\det \gamma_{AB}} \, dy^1 \wedge dy^2$ where $\gamma_{AB}$ is the induced metric on $\Sigma$, and the contraction $\alpha_J(\nu_\gamma) = \frac{1}{|\partial_n|_\gamma}(\alpha_J)_n + (\text{tangential terms})$. The product gives:
\begin{align}
\alpha_J(\nu_\gamma) \, d\sigma_\gamma &= \left(\frac{(\alpha_J)_n}{|\partial_n|_\gamma} + O(\text{tan})\right) \cdot |\partial_n|_\gamma \sqrt{\det \gamma_{AB}} \, dy^1 \wedge dy^2 \\
&= (\alpha_J)_n \sqrt{\det \gamma_{AB}} \, dy^1 \wedge dy^2 + (\text{tangential terms}).
\end{align}

\textit{Step (iii): The tangential terms vanish upon integration.} When we integrate over $\Sigma$, terms involving $\alpha_J(\partial_{y^A})$ for tangent vectors $\partial_{y^A}$ contribute to the boundary $\partial \Sigma$. For closed surfaces ($\partial \Sigma = \emptyset$), these vanish.

\textit{Step (iv): The normal component is metric-independent.} The quantity $(\alpha_J)_n = \alpha_J(\partial_n)$ depends only on the 1-form $\alpha_J$ and the transverse coordinate $n$, not on the metric $\gamma$. The remaining factor $\sqrt{\det \gamma_{AB}}$ appears to depend on $\gamma$, but this is compensated by the implicit dependence of $(\alpha_J)_n$ on the normalization.

More precisely, define the \textbf{metric-free flux 2-form}:
\[
\Phi_{\alpha_J} := \iota^*(\star_g \alpha_J),
\]
where $\star_g$ is the Hodge star with respect to the \textbf{physical} metric $g$. This is a well-defined 2-form on $\Sigma$ depending only on $\alpha_J$, $g$, and the embedding $\iota$. A direct calculation in coordinates shows:
\[
\int_\Sigma \alpha_J(\nu_\gamma) \, d\sigma_\gamma = \int_\Sigma \Phi_{\alpha_J}
\]
for \textbf{any} choice of $\gamma$. The right-hand side is manifestly metric-independent.
\end{proof}

\textbf{Application to the AMO flow.} The level sets $\Sigma_t = \{u = t\}$ are well-defined submanifolds of $M$. We may use $\tg = \phi^4 \bg$ to define their unit normal $\nu_{\tg}$ and area element $d\sigma_{\tg}$, but by the metric-independence above:
\[
J(t) = \int_{\Sigma_t} \alpha_J(\nu_{\tg}) \, d\sigma_{\tg} = \int_{\Sigma_t} (\star_g \alpha_J)|_{\Sigma_t}.
\]
The conservation of $J(t)$ now follows from the closedness of $\star_g \alpha_J$ (i.e., $d(\star_g\alpha_J) = 0$, equivalently the co-closedness $d^\dagger\alpha_J = 0$), which we prove in Step 5.

The key observation is that the Komar 1-form $\alpha_J = \frac{1}{8\pi} K(\eta, \cdot)^\flat$ is defined on the physical manifold, but we integrate it over surfaces $\Sigma_t$ that are level sets in the conformal picture. This is valid because the underlying smooth manifold $M$ is the same (only the metric changes), the level sets $\Sigma_t \subset M$ are well-defined submanifolds independent of which metric we use, the 1-form $\alpha_J$ and its exterior derivative $d\alpha_J$ are tensorial operations that commute with pullback to any submanifold, and the integral $\int_{\Sigma_t} \star_g \alpha_J$ is computed using the physical metric $g$ for the Hodge dual, making it independent of $\tg$.

The co-closedness $d^\dagger\alpha_J = 0$ (equivalently, $d(\star\alpha_J) = 0$) is established on $(M, g)$ using the physical momentum constraint. Once $\star\alpha_J$ is closed, the integral $\int_{\Sigma_t} \star_g \alpha_J$ depends only on the homology class of $\Sigma_t$---this is a topological statement independent of the ambient metric used to define level sets.

\textbf{Metric-independence of the Komar integral.} We now make explicit which quantities use which metric: $\nu_{\tg} := \nabla_{\tg} u / |\nabla_{\tg} u|_{\tg}$ is the unit normal to $\Sigma_t$ with respect to $\tg$; $d\sigma_{\tg}$ is the area element on $\Sigma_t$ induced by $\tg$; and $\alpha_J := \frac{1}{8\pi} K(\eta, \cdot)^\flat_g$ is the Komar 1-form, using the physical metric $g$ to lower the index. The angular momentum integral is:
\[
J(t) = \int_{\Sigma_t} \iota_{\nu_{\tg}} \alpha_J \, d\sigma_{\tg}.
\]
\textbf{Note that}, by Stokes' theorem, if $d(\star\alpha_J) = 0$ (i.e., $\alpha_J$ is co-closed, $d^\dagger\alpha_J = 0$), then:
\[
\int_{\Sigma_{t_1}} \star\alpha_J  = \int_{\Sigma_{t_2}} \star\alpha_J
\]
for surfaces $\Sigma_{t_1}$ and $\Sigma_{t_2}$ that are homologous. This is because the flux integral of a closed 2-form through a surface is a \textbf{topological invariant} depending only on the homology class of $\Sigma$.

More explicitly, let $W = \{t_1 \leq u \leq t_2\}$ be the region between level sets with $\partial W = \Sigma_{t_2} - \Sigma_{t_1}$. Then:
\[
\int_{\Sigma_{t_2}} \star\alpha_J - \int_{\Sigma_{t_1}} \star\alpha_J = \int_W d(\star\alpha_J) = 0.
\]
This identity holds regardless of the metric structure on $W$.

\textbf{Step 1: Orbit space reduction.}
For an axisymmetric 3-manifold $(\tM, \tg)$ with Killing field $\eta = \partial_\phi$, the orbit space is:
\[
\mathcal{Q} := \tM / U(1) \cong \{(r, z) : r \geq 0\},
\]
a 2-dimensional manifold with boundary (the axis $r = 0$). The metric on $\tM$ takes the form:
\[
\tg = g_{\mathcal{Q}} + \rho^2 d\phi^2,
\]
where $g_{\mathcal{Q}}$ is a metric on $\mathcal{Q}$ and $\rho = \rho(r, z) > 0$ is the orbit radius.

\textbf{Step 2: $p$-Harmonic function on orbit space.}
Since the boundary data ($u = 0$ on $\Sigma$, $u \to 1$ at infinity) is axisymmetric and the equation $\Delta_p u = 0$ respects the symmetry, the solution factors through the orbit space:
\[
u: \tM \to \mathbb{R}, \quad u(r, z, \phi) = \bar{u}(r, z),
\]
where $\bar{u}: \mathcal{Q} \to \mathbb{R}$ satisfies a weighted $p$-Laplace equation on $\mathcal{Q}$.

\textbf{Step 3: Gradient orthogonality.}
The gradient of $u$ is:
\[
\nabla u = \nabla_{\mathcal{Q}} \bar{u} + 0 \cdot \partial_\phi,
\]
hence $\nabla u$ lies entirely in $T\mathcal{Q} \subset T\tM$. Since $\eta = \partial_\phi \in T(\text{orbit})$ is orthogonal to $T\mathcal{Q}$:
\[
\tg(\nabla u, \eta) = 0 \quad \text{everywhere on } \tM.
\]
Therefore, the outward unit normal to level sets satisfies:
\[
\nu := \frac{\nabla u}{|\nabla u|} \perp \eta.
\]

\textbf{Step 4: Komar integral as closed form.}
The Komar angular momentum on a surface $\Sigma_t = \{u = t\}$ is:
\[
J(t) = \frac{1}{8\pi} \int_{\Sigma_t} K(\eta, \nu) \, d\sigma = \int_{\Sigma_t} \star_g \alpha_J,
\]
where $\star_g \alpha_J$ is the Hodge dual of the Komar 1-form (a 2-form). For axisymmetric data with $\nu \perp \eta$, Stokes' theorem applied to the 2-form $\star_g \alpha_J$ (or equivalently, via the identity $d(\star \alpha) = \star(d^\dagger \alpha)$ when $\alpha$ is co-closed) yields:
\[
J(t_2) - J(t_1) = \int_{\Sigma_{t_2}} \star_g\alpha_J - \int_{\Sigma_{t_1}} \star_g\alpha_J = \int_{\{t_1 < u < t_2\}} d(\star_g\alpha_J).
\]

\textbf{Step 5: Closedness of Komar form---explicit derivation.}
The key calculation uses the momentum constraint and axisymmetry. Define the 1-form:
\[
\alpha_J := \frac{1}{8\pi} K(\eta, \cdot)^\flat = \frac{1}{8\pi} K_{ij}\eta^i dx^j.
\]
The angular momentum on $\Sigma_t$ is $J(t) = \int_{\Sigma_t} \iota_\nu \alpha_J \, d\sigma$ where $\iota_\nu$ denotes contraction with the normal.

We now prove that $d\alpha_J = 0$ for vacuum axisymmetric data. The exterior derivative of $\alpha_J$ is:
\[
d\alpha_J = \frac{1}{8\pi} d(K_{ij}\eta^i dx^j) = \frac{1}{8\pi} \partial_k(K_{ij}\eta^i) dx^k \wedge dx^j.
\]
Using the product rule:
\begin{equation}\label{eq:dalpha-expansion}
(d\alpha_J)_{kj} = \frac{1}{8\pi}\left[(\nabla_k K_{ij})\eta^i + K_{ij}(\nabla_k \eta^i) - (\nabla_j K_{ik})\eta^i - K_{ik}(\nabla_j \eta^i)\right].
\end{equation}

\textbf{Consolidated proof of co-closedness ($d^\dagger \alpha_J = 0$).}
We now provide a self-contained derivation showing that the Komar 1-form $\alpha_J$ is co-closed for vacuum axisymmetric data, which is the key property ensuring conservation of $J$ via Stokes' theorem.

\textit{Setup.} Define $\beta := K(\eta, \cdot)^\flat$, so $\beta_j = K_{ij}\eta^i$ and $\alpha_J = \frac{1}{8\pi}\beta$. The co-closedness $d^\dagger \alpha_J = 0$ is equivalent to $\nabla^j \beta_j = 0$.

\textit{Computation of $\nabla^j \beta_j$.} Expanding the divergence:
\begin{align}
\nabla^j \beta_j &= \nabla^j (K_{ij}\eta^i) = (\nabla^j K_{ij})\eta^i + K_{ij}(\nabla^j \eta^i). \label{eq:div-beta}
\end{align}

\textit{First term: Momentum constraint.} The momentum constraint reads:
\[
\nabla^j K_{ij} - \nabla_i (\tr K) = 8\pi \momdens_i,
\]
where $\momdens_i$ is the momentum density. Contracting with $\eta^i$:
\[
(\nabla^j K_{ij})\eta^i = 8\pi \momdens_i \eta^i + \eta^i \nabla_i(\tr K) = 8\pi (\momdens \cdot \eta) + \mathcal{L}_\eta(\tr K).
\]
By axisymmetry, $\mathcal{L}_\eta(\tr K) = 0$, so the first term equals $8\pi(\momdens \cdot \eta)$.

\textit{Second term: Killing symmetry.} Using the Killing equation $\nabla^j \eta^i = -\nabla^i \eta^j$:
\[
K_{ij}(\nabla^j \eta^i) = -K_{ij}\nabla^i \eta^j.
\]
Since $K_{ij}$ is symmetric and $\nabla^i \eta^j$ is antisymmetric (Killing equation), their contraction vanishes:
\[
K_{ij}(\nabla^j \eta^i) = 0.
\]

\textit{Conclusion.} Combining these results in \eqref{eq:div-beta}:
\[
\nabla^j \beta_j = 8\pi(j \cdot \eta) + 0 = 8\pi(j \cdot \eta).
\]
Therefore $d^\dagger \alpha_J = \frac{1}{8\pi}\nabla^j \beta_j = j \cdot \eta$. \textbf{For vacuum data ($j = 0$), we obtain $d^\dagger \alpha_J = 0$ exactly.}

\textit{Implication for conservation.} In 3 dimensions, $d^\dagger \alpha_J = 0$ is equivalent to $d(\star_g \alpha_J) = 0$. By Stokes' theorem, for any two homologous surfaces $\Sigma_{t_1}, \Sigma_{t_2}$ bounding region $W$:
\[
J(t_2) - J(t_1) = \int_{\Sigma_{t_2}} \star_g \alpha_J - \int_{\Sigma_{t_1}} \star_g \alpha_J = \int_W d(\star_g \alpha_J) = 0.
\]
This completes the proof that $J(t)$ is constant along the AMO flow for vacuum axisymmetric data.

\begin{remark}[Closedness vs.~co-closedness]\label{rem:closed-vs-coclosed}
The Komar 1-form satisfies $d^\dagger \alpha_J = 0$ (co-closedness), not $d\alpha_J = 0$ (closedness). In 3D, the Hodge dual converts co-closedness of a 1-form to closedness of the corresponding 2-form: $d(\star \alpha) = \star(d^\dagger \alpha)$. Thus $d^\dagger \alpha_J = 0$ implies $d(\star_g \alpha_J) = 0$, which is the condition needed for Stokes' theorem. The distinction matters: $d\alpha_J$ involves derivatives of $K$, while $d^\dagger \alpha_J$ involves the divergence, directly related to the momentum constraint.
\end{remark}

\textit{(Legacy notation---exterior derivative analysis).} For completeness, we record that for vacuum axisymmetric data, $d\beta = 0$ as well. The full exterior derivative $(d\beta)_{jk}$ vanishes because (i) the Killing terms vanish by $\mathcal{L}_\eta K = 0$, and (ii) the momentum constraint terms vanish for $j = 0$. Thus $\alpha_J$ is \textit{both} closed and co-closed for vacuum axisymmetric data, though only co-closedness is needed for the Stokes argument.

\textbf{Step 6: Axisymmetric momentum density.}
For axisymmetric matter satisfying DEC, the momentum density $\momdens_i$ is itself axisymmetric: $\mathcal{L}_\eta \momdens = 0$. On the orbit space $\mathcal{Q} = M/U(1)$, the 1-form $\momdens$ decomposes as $\momdens = \momdens_{\mathcal{Q}} + \momdens_\phi d\phi$. Axisymmetry requires $\momdens_\phi = \momdens_\phi(r, z)$ independent of $\phi$.

The key observation: $\momdens_i\eta^i = \momdens_\phi \cdot |\eta|^2 = \momdens_\phi \rho^2$. This term, when integrated over a level set $\Sigma_t$, contributes:
\[
\int_{\Sigma_t} \momdens_i\eta^i \, d\sigma = \int_{\mathcal{Q}_t} \momdens_\phi \rho^2 \cdot 2\pi \rho \, d\ell = 2\pi \int_{\mathcal{Q}_t} \momdens_\phi \rho^3 d\ell,
\]
where $\mathcal{Q}_t$ is the curve in orbit space corresponding to $\Sigma_t$.

For \textbf{vacuum} data ($\momdens_i = 0$), we have $d\alpha_J = 0$ exactly. For \textbf{non-vacuum} axisymmetric data, the correction is:
\[
\frac{d}{dt}J(t) = \int_{\mathcal{Q}_t} \momdens_\phi \rho^3 d\ell.
\]
Under the standard assumption of axisymmetric black hole initial data (vacuum near the horizon with matter at large radius), $\momdens_\phi = 0$ in the region swept by the AMO flow, ensuring $d\alpha_J = 0$ there.

\textbf{Step 7: Conservation.}
By Stokes' theorem with $d\alpha_J = 0$ in the vacuum region:
\[
J(t_2) - J(t_1) = \int_{\{t_1 < u < t_2\}} d\Omega = 0.
\]
Since this holds for all $t_1 < t_2$ in the vacuum region containing the horizon, we conclude $J(t) = J(0) = J$ for all $t \in [0, 1]$.
\end{proof}

\end{document}
